\section{Data}
\label{sec:data}

DREAM is evaluated on eight real public data sources. We deliberately avoid
any credentialed feed (FAA SWIM, OpenSky Trino) so that the entire pipeline
is reproducible by any researcher with a standard Python installation. All
sources are listed below; \cref{tab:data} reports the per-day ADS-B inventory.

\paragraph{Airports.}
We study two large US hub airports: Los Angeles International (LAX,
ICAO~\texttt{KLAX}) as primary case study, and San Francisco International
(SFO, ICAO~\texttt{KSFO}) as external validation. Both have parallel-runway
configurations whose active direction shifts with the prevailing wind;
this is exactly the situation in which $\Cdyn$ should differ markedly from
$\mathcal{S}_{\mathrm{rwy}}$. Per-airport runway thresholds, ARP, and field
elevation are taken from the FAA NASR-effective Airport Diagrams and
encoded in YAML under \texttt{configs/airports/}.

\paragraph{Annex 14 parameters.}
The OLS dimensions used in \cref{eq:sdf,eq:Astatic} are the
publicly-documented Code~4 precision-approach values, listed in full in
\cref{tab:ols-params}. They drive \emph{both} airports identically; we make
no per-airport adjustment to these parameters.

\paragraph{Operational data: ADS-B from \texttt{adsb.lol}.}
The \texttt{adsb.lol} project~\citep{adsblol_history} publishes a daily
historical archive of \texttt{readsb}-decoded aircraft state vectors,
released as gzipped JSON traces on GitHub. We download the five Friday
releases for August 2024 -- 2024-08-02, 09, 16, 23, 30 -- and filter each
to the 30~NM ARP-centred bounding box of the target airport. We project
WGS-84 coordinates to ENU using the per-airport
azimuthal-equidistant projection. \Cref{tab:data} reports the per-day,
per-airport row counts (\textbf{$7.6$ million state vectors in total}),
the number of distinct aircraft, and the resulting parquet size.

% --- Table 2: per-day data inventory ------------------------------------
\begin{table}[t]
\centering
\caption{Real-data ADS-B inventory per Friday in August 2024, acquired from the open \texttt{adsb.lol} historical archive. Within a 30 NM ARP-centred bounding box. Note the $\sim$50\% coverage drop after 2024-08-09 at both airports, reflecting fewer volunteer receivers on the later Fridays.}
\label{tab:data}
\begin{tabular}{lrrrrrr}
\toprule
 & \multicolumn{3}{c}{LAX (KLAX)} & \multicolumn{3}{c}{SFO (KSFO)} \\
\cmidrule(lr){2-4}\cmidrule(lr){5-7}
Date (Fri) & rows & a/c & MB & rows & a/c & MB \\
\midrule
2024-08-02 & 1,124,625 & 3,763 & 31.0 & 851,836 & 2,069 & 21.8 \\
2024-08-09 & 1,144,920 & 3,851 & 31.8 & 826,573 & 2,084 & 21.1 \\
2024-08-16 & 565,660 & 1,834 & 16.1 & 404,611 & 1,004 & 11.4 \\
2024-08-23 & 498,990 & 1,632 & 14.5 & 366,808 & 899 & 10.4 \\
2024-08-30 & 456,486 & 1,523 & 13.4 & 342,184 & 865 & 9.9 \\
\midrule
\textbf{Sum} & \textbf{3,790,681} & & & \textbf{2,792,012} & & \\
\bottomrule
\end{tabular}
\end{table}


\paragraph{Coverage caveat.}
A noticeable feature of \cref{tab:data} is that the row count drops
$\sim 50\,\%$ at both airports after 2024-08-09. The cause is not an artefact
of our pipeline: \texttt{adsb.lol} is a volunteer-receiver network, and
volunteer-feeder participation on 16/23/30 was lower than on 02/09. This is
exactly the kind of operational coverage variation a deployed DREAM system
would have to live with, and we report it openly. The robustness of our
$B_3 \supsetneq B_2$ conclusion to this coverage variation is part of the
empirical contribution of \cref{sec:lax,sec:sfo}.

\paragraph{Other public sources.}
The remaining six sources are:
(i)~\textbf{FAA Digital Obstacle File}~\citep{faa_dof} for the obstacles
within a 50~NM radius of each ARP;
(ii)~\textbf{USGS 3D Elevation Program (3DEP)}~\citep{usgs_3dep} for the DEM
underneath each airport's 60~km extract box;
(iii)~\textbf{OpenStreetMap (Overpass)}~\citep{osm} for buildings, roads, and
amenities within the 30~km ARP box;
(iv)~\textbf{NOAA Aviation Weather Center} for METAR, retrieved from the
Iowa State University ASOS archive~\citep{iastate_asos} because the live
AWC API only serves the current week;
(v)~\textbf{Los Angeles World Airports (LAWA)}~\citep{lawa2024} for
2024 LAX peak-hour ground-trip counts (this is LAX-only; SFO is not a LAWA
airport);
(vi)~\textbf{BTS Airline Origin-and-Destination Survey (DB1B/DB1C)} for the
passenger-itinerary weights used in the accessibility KPIs.

\paragraph{Operational scope.}
All experiments are evaluated on $5 \times 2 \times 12 \times 4 \times 3 =
1\,440$ corridors: 5~Aug-2024 Fridays, 2~airports, 12~directed vertiport
pairs (V1$\rightarrow$V2 through V4$\rightarrow$V3 over the four candidate
vertiports of \cref{sec:method-sdf}), 4~baselines $B_0\ldots B_3$, and
3~peak hours of the day. The 5~Fridays match the LAWA 2024 Traffic
Generation Report's typical-design-day rule
\citep{lawa2024}.
